\input{../tex-inputs/latex-leseansicht-vorspann}

\section[Arthur und Olga Schnitzler an Gustav Schwarzkopf, 26. 5. 1914]{L04485 Arthur und Olga Schnitzler an Gustav Schwarzkopf, 26. 5. 1914}
\nopagebreak\mylabel{L04485v}
\rehead{ }\normalsize\beginnumbering\briefempfaengerindex{Schwarzkopf, Gustav@\textsc{Schwarzkopf, Gustav}!zzzSchnitzler, Olga@\emph{von Olga Schnitzler}!1914-05-261@{26.\,5.\,1914}|(be}\briefempfaengerindex{Schwarzkopf, Gustav@\textsc{Schwarzkopf, Gustav}!zzzSchnitzler, Arthur@\emph{von Arthur Schnitzler}!1914-05-261@{26.\,5.\,1914}|(be}
\toendnotes[C]{\smallbreak\pagebreak[2]}
\correspDesc{Versand  durch Arthur Schnitzler, Olga Schnitzler am 1914-05-26 in Amsterdam
\newline{}Übermittlung  am 1914-05-26 in Amsterdam
\newline{}Erhalt  durch Gustav Schwarzkopf im Zeitraum [27. 5. 1914
                  – 31. 5. 1914?] in Wien}\toendnotes[C]{\smallbreak}
\Standort{DLA, A:Schnitzler, HS.1985.1.1897.}
\physDesc{Bildpostkarte, 349 Zeichen
\newline{}Handschrift Arthur Schnitzler: Bleistift, deutsche Kurrent
\newline{}Handschrift Olga Schnitzler: Bleistift, lateinische Kurrent
\newline{}Versand: Stempel: »\nobreak{}\oindex{Amsterdam@\textbf{Amsterdam}|pwk}Amsterdam, 26. V. 1914, 11–12 N\nobreak{}«.  }\toendnotes[C]{\smallbreak}\pstart{}{\pb}Hrn \textsc{Gustav Schwarzkopf}\pend{}\pstart{}Wien I\oindex{I., Innere Stadt@\textbf{I., Innere Stadt}|pw}\pend{}\pstart{}\textsc{Tiefer Graben 17}\oindex{Wien@\textbf{Wien}!I., Innere Stadt@\textbf{I., Innere Stadt}!Tiefer Graben 17@\textbf{Tiefer Graben 17}|pw}\pend{}{\bigskip}
\pstart
           \noindent{}{\pb}\textcolor{gray}{\textbf{\begin{otherlanguage}{dutch}\label{K_L04485-1v}\edtext{Branding}{\lemma{\textnormal{\emph{Branding}}}\Cendnote{\textnormal{niederländisch: Brandung}}}\label{K_L04485-1}\end{otherlanguage}.}}\pend
           \vspace{1em}
\pstart
           \textsc{Amsterdam\oindex{Amsterdam@\textbf{Amsterdam}|pw}}{ }26. 5. 914.\pend
           \vspace{0.5em}
\pstart
           Seit Samſtag{ }ſind wir hier, haben viel{ }ſchönes geſehen, waren heut in
                  Zandvoort\oindex{Zandvoort@\textbf{Zandvoort}|pw}, woher dieſe wildwogende Karte
                  sta{\geminationm}t, und grüßen Sie herzlich. Samſtag
               gedenken wir nach Köln\oindex{Köln@\textbf{Köln}|pw}, ein paar Tage {\pb}drauf nach Tutzing\oindex{Tutzing@\textbf{Tutzing}|pw}
               zu reiſen. Auf gutes Wiederſehen\pend
           \pstart \label{T_L04485-1v}\edtext{Ihr \spacefill\mbox{Arthur}}{\lemma{\textnormal{\emph{Ihr Arthur}}}\Cendnote{\textnormal{entlang des rechten Randes
                  geschrieben}}}\label{T_L04485-1}\pend{}\selectlanguage{ngerman}\vspace{1em}
\pstart
           \noindent{}{[}hs. O. Schnitzler:{]} Schöne Grüße an Ihren Bruder\pwindex{Schwarzkopf, Max 12.\,6.\,1857 Wien – 14.\,4.\,1928 ebd.@\textsc{Schwarzkopf, Max} (12.\,6.\,1857 Wien – 14.\,4.\,1928 ebd.)|pwv}!\pend
           \pstart Herzlichst Ihre \spacefill\mbox{O.}\pend{}\selectlanguage{ngerman}\endnumbering\briefempfaengerindex{Schwarzkopf, Gustav@\textsc{Schwarzkopf, Gustav}!zzzSchnitzler, Olga@\emph{von Olga Schnitzler}!1914-05-261@{26.\,5.\,1914}|)be}\briefempfaengerindex{Schwarzkopf, Gustav@\textsc{Schwarzkopf, Gustav}!zzzSchnitzler, Arthur@\emph{von Arthur Schnitzler}!1914-05-261@{26.\,5.\,1914}|)be}\mylabel{L04485h}
\begin{anhang}
\end{anhang}\newcommand{\dateiname}{L04485}\newcommand{\titel}{Arthur und Olga Schnitzler an Gustav Schwarzkopf, 26. 5. 1914}\newcommand{\editorInnen}{Herausgegeben von Jahnke, SelmaMüller, Martin Anton}\input{../tex-inputs/latex-leseansicht-abspann}
